\Needspace{7\baselineskip}\Article{306}
Der Beweis des vorigen Satzes ist den Beweisen in Art.\ \(45,49\) ganz analog. Tatsächlich hat die Theorie der Multiplikation von Klassen in jeder Hinsicht große Ähnlichkeit mit dem in Abschnitt III behandelten Gegenstand. Die Grenzen dieses Werks erlauben aber nicht, diese Theorie hier so ausführlich zu verfolgen, wie sie es verdient. Deshalb fügen wir nur wenige Bemerkungen hinzu, übergehen auch die Beweise, die längere Vorbereitungen erfordern würden, und behalten uns eine umfassendere Untersuchung für eine andere Gelegenheit vor.

I.\ Setzt man die Reihe \(K,C,2C,3C\) usw.\ über \((m-1)C\) hinaus fort, so treten dieselben Klassen wieder auf:
\[mC=K,\quad(m+1)C=C,\quad(m+2)C=2C\text{ usw.}\]
Allgemein sind, wenn man der Übersicht wegen \(K\) als \(0C\) betrachtet, die Klassen \(gC,g'C\) identisch oder verschieden, je nachdem \(g\) und \(g'\) nach dem Modul \(m\) kongruent oder inkongruent sind. Die Klasse \(nC\) ist also stets mit der Hauptklasse \(K\) identisch.

II.\ Die Gesamtheit der Klassen \(K,C,2C\ldots(m-1)C\), die wir oben mit \(\mathfrak C\) bezeichnet haben, nennen wir die \textit{Periode} der Klasse \(C\). Diesen Ausdruck darf man nicht mit den \textit{Perioden reduzierter Formen} positiver nichtquadratischer Determinante verwechseln, die in Art.\ \(186\) und den folgenden Artikeln behandelt wurden. Offenbar entsteht durch Zusammensetzung beliebig vieler Klassen derselben Periode wieder eine Klasse dieser Periode:
\[gC+g'C+g''C\text{ usw.}=(g+g'+g''+\text{usw.})C\]


III.\ Da \(C+(m-1)C=K\), sind die Klassen \(C\) und \((m-1)C\) entgegengesetzt, ebenso \(2C\) und \((m-2)C\), \(3C\) und \((m-3)C\) usw.\ Ist also \(m\) gerade, so ist die Klasse \(\tfrac12mC\) sich selbst entgegengesetzt und damit \textit{ambig}. Kommt umgekehrt in \(\mathfrak C\) außer \(K\) eine weitere ambige Klasse vor, etwa \(gC\), so gilt \(gC=(m-g)C\), also \(g=m-g=\tfrac12m\). Daraus folgt: Ist \(m\) gerade, so kann es außer \(K\) und \(\tfrac12mC\) keine weitere ambige Klasse geben;
\SourceLeaf{377}
ist dagegen \(m\) ungerade, so kann außer \(K\) keine weitere in \(\mathfrak C\) enthalten sein.

IV.\ Die Periode einer Klasse \(hC\), die in \(\mathfrak C\) enthalten ist, sei
\[K,hC,2hC,3hC\ldots(m'-1)hC\]
Dann ist offenbar \(m'h\) das kleinste Vielfache von \(h\), das durch \(m\) teilbar ist. Sind also \(m\) und \(h\) teilerfremd, so ist \(m'=m\), und die beiden Perioden enthalten dieselben Klassen in verschiedener Reihenfolge. Bezeichnet allgemein \(\mu\) den größten gemeinsamen Teiler von \(m,h\), so ist \(m'=\frac m\mu\). Die Anzahl der Klassen in der Periode einer beliebigen Klasse aus \(\mathfrak C\) ist daher entweder \(m\) oder ein echter Teiler von \(m\). Und zwar haben so viele Klassen in \(\mathfrak C\) Perioden mit \(m\) Gliedern, wie unter den Zahlen \(0,1,2\ldots m-1\) zu \(m\) teilerfremd sind, also \(\varphi m\) mit dem Zeichen aus Art.\ \(39\). Allgemein haben so viele Klassen in \(\mathfrak C\) Perioden mit \(\frac m\mu\) Gliedern, wie unter den Zahlen \(0,1,2\ldots m-1\) den größten gemeinsamen Teiler \(\mu\) mit \(m\) haben. Ihre Anzahl ist offenbar \(\varphi\frac m\mu\). Gilt also \(m=n\), das heißt, ist das \textit{ganze} Hauptgeschlecht in \(\mathfrak C\) enthalten, so gibt es in diesem Geschlecht genau \(\varphi n\) Klassen, deren Perioden dasselbe ganze Geschlecht umfassen, und \(\varphi e\) Klassen, deren Perioden aus \(e\) Gliedern bestehen, wobei \(e\) einen beliebigen Teiler von \(n\) bezeichnet. Dieser Schluss gilt allgemein, wenn im Hauptgeschlecht eine Klasse vorkommt, deren Periode aus \(n\) Gliedern besteht.

V.\ Unter derselben Voraussetzung lässt sich das System der Klassen des Hauptgeschlechts am zweckmäßigsten ordnen, indem man eine Klasse mit einer Periode von \(n\) Gliedern gewissermaßen als \textit{Basis} wählt und die Klassen des Hauptgeschlechts in der Reihenfolge ihrer Periode anordnet. Gibt man dann der Hauptklasse den \textit{Index} \(0\), der als Basis gewählten Klasse den Index \(1\) usw., so lässt sich allein durch Addition der Indizes feststellen, welche Klasse aus der Zusammensetzung beliebiger Klassen des Hauptgeschlechts entsteht. Hier ein Beispiel für die Determinante \(-356\), bei dem wir die Klasse \((9,2,40)\) als Basis gewählt haben:
\begin{center}\begin{tabular}{rl|rl|rl}
\(0\) & \((1,0,356)\) & \(4\) & \((20,8,21)\) & \(8\) & \((20,-8,21)\)\\[2pt]
\(1\) & \((9,2,40)\) & \(5\) & \((17,1,21)\) & \(9\) & \((8,2,45)\)\\[2pt]
\(2\) & \((5,2,72)\) & \(6\) & \((4,0,89)\) & \(10\) & \((5,-2,72)\)\\[2pt]
\(3\) & \((8,-2,45)\) & \(7\) & \((17,-1,21)\) & \(11\) & \((9,-2,40)\)\\[2pt]
\end{tabular}\end{center}

\SourceLeaf{378}

VI.\ Sowohl die Analogie mit Abschnitt III als auch die Untersuchung von mehr als \(200\) negativen und noch weit mehr positiven nichtquadratischen Determinanten mögen es höchst wahrscheinlich erscheinen lassen, dass jene Voraussetzung für \textit{alle} Determinanten gilt. Ein solcher Schluss wäre trotzdem falsch und würde durch Fortsetzung der Klassifikationstafel widerlegt. Der Kürze wegen nennen wir die Determinanten, bei denen das ganze Hauptgeschlecht in einer einzigen Periode enthalten sein kann, \textit{regulär}, die übrigen, bei denen das nicht möglich ist, \textit{irregulär}. Dieser Gegenstand scheint zu den tiefsten Geheimnissen der höheren Arithmetik zu gehören und Raum für äußerst schwierige Untersuchungen zu lassen. Hier können wir ihn nur durch wenige Bemerkungen erläutern, denen wir die folgende allgemeine voranstellen.

VII.\ Gibt es im Hauptgeschlecht Klassen \(C,C'\), deren Perioden aus \(m,m'\) Klassen bestehen, und ist \(M\) die kleinste durch \(m\) und \(m'\) teilbare Zahl, so gibt es in demselben Geschlecht auch Klassen, deren Perioden \(M\) Glieder enthalten. Man zerlege \(M\) in zwei teilerfremde Faktoren \(r,r'\), deren erster \((r)\) die Zahl \(m\), deren zweiter \((r')\) die Zahl \(m'\) teilt, wie in Art.\ \(73\). Dann hat die Klasse \(\frac mrC+\frac{m'}{r'}C'=C''\) die verlangte Eigenschaft. Bestehe nämlich die Periode der Klasse \(C''\) aus \(g\) Gliedern, so gilt
\[\begin{aligned}K=grC''&=gmC+\frac{grm'}{r'}C'\\&=K+\frac{grm'}{r'}C'=\frac{grm'}{r'}C'\end{aligned}\]
also muss \(\frac{grm'}{r'}\) durch \(m'\) teilbar sein oder \(gr\) durch \(r'\) und damit auch \(g\) durch \(r'\). Ganz ebenso findet man \(g\) durch \(r\) teilbar und daher auch durch \(rr'=M\). Da aber offenbar \(MC''=K\), ist auch \(M\) durch \(g\) teilbar. Also gilt notwendig \(M=g\). Daraus folgt leicht, dass die \textit{größte} Anzahl von Klassen in einer Periode bei gegebener Determinante durch die Klassenanzahl jeder anderen Periode einer Klasse desselben Hauptgeschlechts teilbar ist. Zugleich lässt sich daraus eine Methode gewinnen, eine Klasse mit möglichst großer Periode zu finden, die also bei regulärer Determinante das ganze Hauptgeschlecht umfasst. Sie ist der Methode der Art.\ \(73,74\) ganz analog, obwohl sich die Arbeit in der Praxis durch verschiedene Kunstgriffe verkürzen lässt. Der Quotient aus \(n\) und der Klassenanzahl der größten Periode ist für reguläre Determinanten \(1\), für irreguläre aber stets eine ganze Zahl größer als \(1\). Bei diesen eignet er sich besonders gut, die verschiedenen Arten der Irregularität auszudrücken, und kann deshalb \textit{Irregularitätsexponent} heißen.
\SourceLeaf{379}

VIII.\ Bisher gibt es keine allgemeine Regel, mit der sich reguläre von irregulären Determinanten im Voraus unterscheiden lassen, zumal unter den letzten sowohl Primzahlen als auch zusammengesetzte Zahlen vorkommen. Einige besondere Beobachtungen müssen hier also genügen. Enthält das Hauptgeschlecht mehr als zwei ambige Klassen, so ist die Determinante sicher irregulär und der Irregularitätsexponent gerade. Gibt es in diesem Geschlecht dagegen nur eine oder zwei, so ist die Determinante entweder regulär oder wenigstens ihr Irregularitätsexponent ungerade. Alle negativen Determinanten der Form \(-(216k+27)\), mit der einzigen Ausnahme \(-27\), sind irregulär, und ihr Irregularitätsexponent ist durch \(3\) teilbar. Dasselbe gilt für negative Determinanten der Formen \(-(1000k+75)\) und \(-(1000k+675)\), mit der einzigen Ausnahme \(-75\), und für unendlich viele andere. Ist der Irregularitätsexponent eine Primzahl \(p\) oder wenigstens durch \(p\) teilbar, so ist \(n\) durch \(pp\) teilbar. Hat also \(n\) keinen quadratischen Teiler, so ist die Determinante sicher regulär. Nur bei positiven Determinanten, die \textit{Quadratzahlen} \(ee\) sind, lässt sich stets im Voraus erkennen, ob sie regulär oder irregulär sind. Regulär sind sie, wenn \(e\) gleich \(1\) oder \(2\), eine ungerade Primzahl oder eine Potenz einer ungeraden Primzahl ist; in allen übrigen Fällen sind sie irregulär. Bei negativen Determinanten werden die irregulären mit wachsender Determinante immer häufiger. \textit{Zum Beispiel} findet man im ganzen ersten Tausender dreizehn irreguläre: ohne negatives Vorzeichen \(576,580,820,884,900\) mit dem Irregularitätsexponenten \(2\) sowie \(243,307,339,459,675,755,891,974\) mit dem Irregularitätsexponenten \(3\). Im zweiten Tausender wurden \(13\) mit dem Exponenten \(2\) und \(15\) mit dem Exponenten \(3\) gefunden; im zehnten Tausender \(31\) mit dem Exponenten \(2\) und \(32\) mit dem Exponenten \(3\). Ob Determinanten mit einem Irregularitätsexponenten größer als \(3\) unter \(-10000\) vorkommen, lässt sich noch nicht entscheiden. Jenseits dieser Grenze können beliebige gegebene Exponenten vorkommen. Es ist sehr wahrscheinlich, dass mit wachsenden Determinanten das Verhältnis der Häufigkeit irregulärer negativer Determinanten zu der regulärer sich immer mehr einer festen Zahl nähert. Diese zu bestimmen wäre des Scharfsinns der Mathematiker sehr würdig. — Bei positiven nichtquadratischen Determinanten sind irreguläre viel seltener. Mit geradem Irregularitätsexponenten gibt es sicher unendlich viele, \textit{zum Beispiel} \(3026\), bei der er \(2\) ist. Es scheint auch kein Zweifel, dass solche mit ungeradem Irregularitätsexponenten vorkommen, obwohl wir zugeben müssen, dass uns bisher keine begegnet ist.

IX.\ Die zweckmäßigste Anordnung des Systems der Klassen im Hauptgeschlecht einer irregulären Determinante können wir hier der Kürze wegen nicht behandeln. Wir bemerken nur: Da eine einzige Basis hier nicht genügt, muss man zwei
\SourceLeaf{380}
oder sogar noch mehr Klassen wählen, aus deren Multiplikation und Zusammensetzung alle Klassen entstehen. Daraus ergeben sich \textit{doppelte oder mehrfache Indizes}, die ungefähr denselben Zweck erfüllen wie die einfachen bei regulären Determinanten. Dies behandeln wir ein anderes Mal ausführlicher.

X.\ Schließlich bemerken wir, dass alle in diesem und dem vorigen Artikel betrachteten Eigenschaften vor allem von der Zahl \(n\) abhängen, die eine ähnliche Rolle spielt wie \(p-1\) in Abschnitt III. Daher verdient diese Zahl die größte Aufmerksamkeit, und eine allgemeine Beziehung zwischen ihr und der zugehörigen Determinante zu entdecken wäre höchst wünschenswert. Bei dieser sehr wichtigen Frage halten wir es umso weniger für angebracht, die Hoffnung aufzugeben, als es bereits gelungen ist, den Durchschnittswert des Produkts aus \(n\) und der im Voraus bestimmbaren Geschlechteranzahl wenigstens bei negativen Determinanten durch eine analytische Formel auszudrücken, in Art.\ \(302\).

\Needspace{7\baselineskip}\Article{307}
Die Untersuchungen der vorigen Artikel umfassen nur die Klassen des Hauptgeschlechts. Sie genügen daher sowohl für positive Determinanten mit nur einem Geschlecht überhaupt als auch für negative mit nur einem positiven Geschlecht, wenn wir das negative nicht berücksichtigen wollen. Es bleibt, auch über die übrigen eigentlich primitiven Geschlechter einige Bemerkungen hinzuzufügen.

I.\ Gibt es in einem Geschlecht \(G'\), das vom Hauptgeschlecht \(G\) derselben Determinante verschieden ist, eine ambige Klasse, so gibt es darin ebenso viele wie in \(G\). Die ambigen Klassen in \(G\) seien \(L\), \(M,N\) usw., darunter auch die Hauptklasse \(K\), in \(G'\) dagegen \(L',M',N'\) usw.\ Ihre Gesamtheiten seien der Reihe nach \(A\) und \(A'\). Da offenbar alle Klassen \(L+L',\ M+L',\ N+L'\) usw.\ ambig und verschieden sind, zu \(G'\) gehören und somit in \(A'\) enthalten sein müssen, kann die Klassenanzahl in \(A'\) sicher nicht kleiner sein als in \(A\). Ebenso sind die Klassen \(L'+L'\), \(M'+L',\ N'+L'\) usw.\ verschieden und ambig, gehören zu \(G\) und sind also in \(A\) enthalten. Daher kann die Klassenanzahl in \(A\) nicht kleiner sein als in \(A'\). Die Klassenanzahlen in \(A\) und \(A'\) sind somit notwendig gleich.

II.\ Da die Gesamtzahl der ambigen Klassen gleich der Anzahl der Geschlechter ist, nach Art.\ \(261,287\) III, muss offenbar, wenn es in \(G\) nur eine ambige
\SourceLeaf{381}
Klasse gibt, \textit{jedes} Geschlecht eine ambige Klasse enthalten. Gibt es in \(G\) zwei, so enthält die Hälfte aller Geschlechter je zwei, die übrigen keine. Sind schließlich in \(G\) mehr enthalten, etwa \(a\)\HistoricalNote{\ensuremath{\ast})}{Dies kann nur bei irregulären Determinanten geschehen, und \(a\) ist dann stets eine Zweierpotenz.}, so enthält der \(a\)-te Teil aller Geschlechter je \(a\) ambige Klassen, die übrigen keine.

III.\ Für den Fall, dass \(G\) zwei ambige Klassen enthält, seien \(G,G',G''\) usw.\ die Geschlechter mit je zwei und \(H,H',H''\) usw.\ die ohne ambige Klassen. Die Gesamtheit der ersten sei \(\mathfrak G\), die der letzten \(\mathfrak H\). Da durch Zusammensetzung zweier ambiger Klassen stets eine ambige Klasse entsteht, nach Art.\ \(249\), ergibt sich durch Zusammensetzung zweier Geschlechter aus \(\mathfrak G\) offenbar stets ein Geschlecht aus \(\mathfrak G\). Daraus folgt weiter, dass die Zusammensetzung eines Geschlechts aus \(\mathfrak G\) mit einem aus \(\mathfrak H\) ein Geschlecht aus \(\mathfrak H\) ergibt. Denn gehörte \textit{zum Beispiel} \(G'+H\) nicht zu \(\mathfrak H\), sondern zu \(\mathfrak G\), so müsste auch \(G'+H+G'\) zu \(\mathfrak G\) gehören. \textit{Das ist ein Widerspruch}, da \(G'+G'=G\) und somit \(G'+H+G'=H\). Schließlich sieht man leicht, dass die Geschlechter \(G+H\), \(G'+H,\ G''+H\) usw.\ zusammen mit \(H+H,\ H'+H,\ H''+H\) usw.\ alle verschieden und daher mit \(\mathfrak G\) und \(\mathfrak H\) zusammengenommen identisch sind. Nach dem gerade Bewiesenen gehören aber die Geschlechter \(G+H,\ G'+H,\ G''+H\) usw.\ alle zu \(\mathfrak H\) und schöpfen diese Gesamtheit also aus. Daher müssen die übrigen \(H+H,\ H'+H\), \(H''+H\) usw.\ sämtlich zu \(\mathfrak G\) gehören. \textit{Das heißt}, die Zusammensetzung zweier Geschlechter aus \(\mathfrak H\) ergibt also stets ein Geschlecht aus \(\mathfrak G\).

IV.\ Ist \(E\) eine Klasse des Geschlechts \(V\), das vom Hauptgeschlecht \(G\) verschieden ist, so gehören offenbar \(2E,4E\), \(6E\) usw.\ alle zu \(G\), dagegen \(3E,5E,7E\) usw.\ zu \(V\). Besteht die Periode der Klasse \(2E\) aus \(m\) Gliedern, so ist in der Reihe \(E,2E,3E\) usw.\ offenbar die Klasse \(2mE\) und keine frühere mit \(K\) identisch. Die Periode der Klasse \(E\) besteht also aus \(2m\) Gliedern. Daher ist die Gliederzahl der Periode einer beliebigen Klasse außerhalb des Hauptgeschlechts entweder \(2n\) oder ein echter Teiler von \(2n\), wobei \(n\) die Anzahl der Klassen in jedem Geschlecht bezeichnet.

V.\ Es sei \(C\) eine gegebene Klasse des Hauptgeschlechts \(G\) und \(E\) eine Klasse eines Geschlechts \(V\), aus deren Verdopplung \(C\) entsteht; eine solche gibt es stets nach Art.\ \(286\). Ferner seien sämtliche ambigen Klassen, eigentlich primitiv und von derselben Determinante, \(K,K',K''\) usw.\ Dann sind \textit{alle} Klassen, aus deren
\SourceLeaf{382}
Verdopplung \(C\) entsteht, die folgenden: \(E\ (=E+K),\ E+K',\ E+K''\) usw.\ Ihre Gesamtheit sei \(\mathfrak Q\). Die Anzahl dieser Klassen ist gleich der Anzahl der ambigen Klassen, also der Anzahl der Geschlechter. Offenbar gehören von den Klassen in \(\mathfrak Q\) so viele zum Geschlecht \(V\), wie es ambige Klassen in \(G\) gibt. Ihre Anzahl sei \(a\). Dann enthält jedes Geschlecht entweder \(a\) Klassen aus \(\mathfrak Q\) oder keine. Daraus folgt leicht: Ist \(a=1\), so enthält jedes Geschlecht eine Klasse aus \(\mathfrak Q\); ist \(a=2\), so enthält die Hälfte aller Geschlechter je zwei Klassen aus \(\mathfrak Q\), die übrigen keine. Dabei stimmt die erste Hälfte entweder ganz mit \(\mathfrak G\) überein, im selben Sinn wie oben in III, und die zweite mit \(\mathfrak H\), oder diese mit \(\mathfrak G\) und jene mit \(\mathfrak H\). — Ist \(a\) noch größer, so enthält stets der \(a\)-te Teil aller Geschlechter Klassen aus \(\mathfrak Q\), und zwar jedes dieser Geschlechter \(a\) Klassen.

VI.\ Nun sei \(C\) eine Klasse, deren Periode aus \(n\) Gliedern besteht. Für den Fall \(a=2\), in dem also \(n\) gerade ist, sieht man leicht, dass keine Klasse aus \(\mathfrak Q\) zu \(G\) gehören kann. Dann wäre nämlich eine solche Klasse in der Periode der Klasse \(C\) enthalten. Wäre sie also \(=rC\), oder \(2rC=C\), so wäre \(2r\equiv1\; (\mathrm{mod}.\ n)\), \textit{was ein Widerspruch ist}. Da \(G\) zu \(\mathfrak G\) gehört, müssen also sämtliche Klassen \(\mathfrak Q\) auf die Geschlechter \(\mathfrak H\) verteilt sein. Weil es bei regulärer Determinante in \(G\) genau \(\varphi n\) Klassen mit Perioden von \(n\) Gliedern gibt, folgt: Für den Fall \(a=2\) enthält jedes Geschlecht \(\mathfrak H\) genau \(2\varphi n\) Klassen, deren Perioden \(2n\) Glieder und somit sowohl ihr eigenes als auch das Hauptgeschlecht umfassen. Ist dagegen \(a=1\), so gibt es in jedem vom Hauptgeschlecht verschiedenen Geschlecht \(\varphi n\) solche Klassen.

VII.\ Auf diese Beobachtungen gründen wir die folgende Methode, das System \textit{aller} eigentlich primitiven Klassen für jede gegebene reguläre Determinante möglichst zweckmäßig aufzubauen; irreguläre lassen wir ganz beiseite. Man wähle beliebig eine Klasse \(E\), deren Periode \(2n\) Glieder und damit sowohl ihr eigenes Geschlecht, etwa \(V\), als auch das Hauptgeschlecht \(G\) umfasst. Die Klassen dieser beiden Geschlechter ordne man in der Reihenfolge dieser Periode. Damit ist die Aufgabe schon erledigt, wenn es überhaupt nicht mehr als diese beiden Geschlechter gibt oder wenn es nicht nötig erscheint, die übrigen hinzuzufügen, \textit{zum Beispiel} bei einer negativen Determinante mit nur zwei positiven Geschlechtern. Sollen aber vier oder mehr Geschlechter aufgebaut werden, so behandle man die übrigen auf folgende Weise. Es sei \(V'\) eines der übrigen und \(V+V'=V''\). Dann gibt es in \(V'\) und \(V''\) zusammen zwei ambige Klassen, entweder in jedem eine oder in einem zwei und im anderen keine. Daraus
\SourceLeaf{383}
wähle man eine beliebige \(A\). Setzt man \(A\) mit den einzelnen Klassen in \(G\) und \(V\) zusammen, so entstehen offenbar \(2n\) verschiedene Klassen, die zu \(V'\) und \(V''\) gehören und diese Geschlechter somit ausschöpfen. So lassen sich auch diese Geschlechter ordnen. — Sind außer diesen vier weitere Geschlechter übrig, sei \(V'''\) eines davon, und es seien \(V'''',V''''',V''''''\) die Geschlechter, die durch Zusammensetzung von \(V'''\) mit \(V,V'\) und \(V''\) entstehen. Diese vier Geschlechter \(V'''\ldots V''''''\) enthalten vier ambige Klassen. Wählt man daraus eine \(A'\) und setzt sie mit den einzelnen Klassen in \(G,V,V',V''\) zusammen, so entstehen offenbar alle Klassen in \(V'''\ldots V''''''\). — Sind noch weitere Geschlechter übrig, so fahre man ebenso fort, bis alle ausgeschöpft sind. Ist die Anzahl aller aufzubauenden Geschlechter \(2^\mu\), so braucht man offenbar insgesamt \(\mu-1\) ambige Klassen. Jede Klasse dieser Geschlechter entsteht entweder durch Multiplikation der Klasse \(E\) oder durch Zusammensetzung einer so entstandenen Klasse mit einer oder mehreren ambigen Klassen. Zwei Beispiele erläutern diese Regeln. Mehr über den Nutzen einer solchen Anordnung oder über Kunstgriffe zur Erleichterung der Arbeit können wir hier nicht hinzufügen.

\Needspace{12\baselineskip}\begin{center}I. Determinante \(-161\).\par\smallskip Vier positive Geschlechter; je vier Klassen.\end{center}
\Needspace{9\baselineskip}
\begin{center}
\begin{minipage}[t]{0.47\linewidth}\centering
\(G\)\par\smallskip
{\small \(1,4;\ R7;\ R23\)}\par\smallskip
\begin{tabular}{@{}l@{}}
\((1,0,161)=K\)\\\((9,1,18)=2E\)\\\((2,1,81)=4E\)\\\((9,-1,18)=6E\)
\end{tabular}
\end{minipage}\hfill\begin{minipage}[t]{0.47\linewidth}\centering
\(V\)\par\smallskip
{\small \(3,4;\ N7;\ R23\)}\par\smallskip
\begin{tabular}{@{}l@{}}
\((3,1,54)=E\)\\\((6,-1,27)=3E\)\\\((6,1,27)=5E\)\\\((3,-1,54)=7E\)
\end{tabular}
\end{minipage}
\end{center}
\Needspace{9\baselineskip}
\begin{center}
\begin{minipage}[t]{0.47\linewidth}\centering
\(V'\)\par\smallskip
{\small \(3,4;\ R7;\ N23\)}\par\smallskip
\begin{tabular}{@{}l@{}}
\((7,0,23)=A\)\\\((11,-2,15)=A+2E\)\\\((14,7,15)=A+4E\)\\\((11,2,15)=A+6E\)
\end{tabular}
\end{minipage}\hfill\begin{minipage}[t]{0.47\linewidth}\centering
\(V''\)\par\smallskip
{\small \(1,4;\ N7;\ N23\)}\par\smallskip
\begin{tabular}{@{}l@{}}
\((10,3,17)=A+E\)\\\((5,2,33)=A+3E\)\\\((5,-2,33)=A+5E\)\\\((10,-3,17)=A+7E\)
\end{tabular}
\end{minipage}
\end{center}
\SourceLeaf{384}
\Needspace{12\baselineskip}\begin{center}II. Determinante \(-546\).\par\smallskip Acht positive Geschlechter; je drei Klassen.\end{center}
\Needspace{9\baselineskip}
\begin{center}
\begin{minipage}[t]{0.47\linewidth}\centering
\(G\)\par\smallskip
{\small \(1\text{ und }3, 8; R3; R7; R13\)}\par\smallskip
\begin{tabular}{@{}l@{}}
\((1,0,546)=K\)\\\((22,-2,25)=2E\)\\\((22,2,25)=4E\)
\end{tabular}
\end{minipage}\hfill\begin{minipage}[t]{0.47\linewidth}\centering
\(V\)\par\smallskip
{\small \(5\text{ und }7, 8; N3; N7; N13\)}\par\smallskip
\begin{tabular}{@{}l@{}}
\((5,2,110)=E\)\\\((21,0,26)=3E\)\\\((5,-2,110)=5E\)
\end{tabular}
\end{minipage}
\end{center}
\Needspace{9\baselineskip}
\begin{center}
\begin{minipage}[t]{0.47\linewidth}\centering
\(V'\)\par\smallskip
{\small \(1\text{ und }3, 8; N3; R7; N13\)}\par\smallskip
\begin{tabular}{@{}l@{}}
\((2,0,273)=A\)\\\((11,-2,50)=A+2E\)\\\((11,2,50)=A+4E\)
\end{tabular}
\end{minipage}\hfill\begin{minipage}[t]{0.47\linewidth}\centering
\(V''\)\par\smallskip
{\small \(5\text{ und }7, 8; R3; N7; R13\)}\par\smallskip
\begin{tabular}{@{}l@{}}
\((10,2,55)=A+E\)\\\((13,0,42)=A+3E\)\\\((10,-2,55)=A+5E\)
\end{tabular}
\end{minipage}
\end{center}
\Needspace{9\baselineskip}
\begin{center}
\begin{minipage}[t]{0.47\linewidth}\centering
\(V'''\)\par\smallskip
{\small \(1\text{ und }3, 8; N3; N7; R13\)}\par\smallskip
\begin{tabular}{@{}l@{}}
\((3,0,182)=A'\)\\\((17,7,35)=A'+2E\)\\\((17,-7,35)=A'+4E\)
\end{tabular}
\end{minipage}\hfill\begin{minipage}[t]{0.47\linewidth}\centering
\(V''''\)\par\smallskip
{\small \(5\text{ und }7, 8; R3; R7; N13\)}\par\smallskip
\begin{tabular}{@{}l@{}}
\((15,-3,37)=A'+E\)\\\((7,0,78)=A'+3E\)\\\((15,3,37)=A'+5E\)
\end{tabular}
\end{minipage}
\end{center}
\Needspace{9\baselineskip}
\begin{center}
\begin{minipage}[t]{0.47\linewidth}\centering
\(V'''''\)\par\smallskip
{\small \(1\text{ und }3, 8; R3; N7; N13\)}\par\smallskip
\begin{tabular}{@{}l@{}}
\((6,0,91)=A+A'\)\\\((19,9,33)=A+A'+2E\)\\\((19,-9,33)=A+A'+4E\)
\end{tabular}
\end{minipage}\hfill\begin{minipage}[t]{0.47\linewidth}\centering
\(V''''''\)\par\smallskip
{\small \(5\text{ und }7, 8; N3; R7; R13\)}\par\smallskip
\begin{tabular}{@{}l@{}}
\((23,11,29)=A+A'+E\)\\\((14,0,39)=A+A'+3E\)\\\((23,-11,29)=A+A'+5E\)
\end{tabular}
\end{minipage}
\end{center}\BookRule{0.2\linewidth}
